% Build from paper/: pdflatex paper && bibtex paper && pdflatex paper && pdflatex paper
\documentclass[conference]{IEEEtran}
\usepackage[T1]{fontenc}
\usepackage[utf8]{inputenc}
\usepackage{lmodern}
\usepackage[expansion=false]{microtype}
\usepackage{amsmath,amssymb}
\usepackage{booktabs,tabularx,array}
\usepackage{xcolor}
\usepackage{listings}
\usepackage{cite}
\usepackage[hidelinks]{hyperref}
\urlstyle{same}
\hypersetup{pdftitle={Kagura: Policy-Driven LLVM Obfuscation for Native Clients},
  pdfauthor={Yoshiki Kusama}}
\lstset{basicstyle=\ttfamily\footnotesize,breaklines=true,
  columns=fullflexible,keepspaces=true,frame=single,
  keywordstyle=\color{blue!60!black},showstringspaces=false}

\begin{document}
\title{Kagura: Policy-Driven LLVM Obfuscation\\for Native Clients}
\author{\IEEEauthorblockN{Yoshiki Kusama}
\IEEEauthorblockA{Independent Security Researcher}}
\maketitle

\begin{abstract}
Deploying compiler-based obfuscation requires correct pass composition,
consistent configuration, and runtime checks that operate in the target
environment. We present Kagura, an open-source LLVM New Pass Manager
toolkit with a shared pass registry, policy-based selection, and conservative
transformation guards. We examine its core mechanisms and evaluate both
selected pass combinations and the current shipped profiles on LLVM~21.1.8
and arm64 macOS. A nine-pass experiment measures executable text growth
of $1.64$--$1.66\times$ and $1.83$--$1.89\times$ for two explicit
configurations, with seven self-tests passing in every build.
A separate deployment experiment exposes a limitation: all three shipped
profiles abort during startup checks on the evaluation host. With anti-debug
instrumentation explicitly disabled, all nine profile/seed combinations pass
the self-tests and match 4,096 baseline return values each. Subject-module
text growth in that ablation ranges from $1.00\times$ for FAST to
$8.05$--$9.40\times$ for STRONG under its stated compilation protocol.
All 84 registered CTest tests pass, illustrating why component tests must be
complemented by policy-level execution. These results characterize correctness,
code size, repeatability, and deployment outcomes; they do not establish
attack resistance or predict analyst effort.
\end{abstract}

\begin{IEEEkeywords}
code obfuscation, LLVM, compiler passes, game security, reproducibility
\end{IEEEkeywords}

\section{Introduction}
\label{sec:intro}

Native game clients execute in environments that their operators may
inspect and modify. An attacker can search a binary for validation
routines, trace game-state updates, or patch a local check. Keeping
currency, entitlements, and competitive outcomes authoritative on a
server reduces the consequences of client compromise, but native
client code can still contain logic that developers want to make more
costly to inspect or modify. Compiler-level obfuscation is one component
of that engineering problem.

Prior work provides both obfuscating transformations and techniques
for recovering their semantics~\cite{collberg1997taxonomy,junod2015obfuscator,
yadegari2015generic,blazytko2017syntia}. Consequently, the presence of
flattened control flow or mixed Boolean-arithmetic (MBA) expressions
is not itself evidence of security. An implementation must first
preserve the behavior of the program, and its costs and security
properties must be measured separately.

Kagura implements obfuscation and runtime support around LLVM's New
Pass Manager (NPM). This paper examines the engineering of a composable compiler pipeline,
using nine selected passes to explain its mechanisms, followed by a separate
evaluation of the current shipped profiles. The project contains a broader
catalogue of transformations and platform integrations; neither experiment
exercises every transformation or platform. The algorithms in the selected
passes are established techniques. The contribution is their
implementation and composition in a maintained toolkit, together with
an evaluation protocol that makes its scope and evidence explicit.

We address three questions:
\begin{enumerate}
  \item How are pass selection, ordering, and semantic safeguards
        represented in the implementation?
  \item What native code-size costs do two explicitly defined pass
        combinations incur on small game-related verification routines?
  \item Do the shipped policies execute on the evaluation host, and how do
        those outcomes differ from component tests and selected pass lists?
\end{enumerate}

The resulting artifact includes the pass implementations, a native
measurement runner, raw results, and generated paper tables. It is
intended as a basis for repeatable experiments, rather than a claim
that a particular number of transformations achieves a particular
level of protection.

\section{Scope and Threat Model}
\label{sec:threat}

\subsection{Assets and Adversary}

The assets of interest are the structure of client-side verification
logic and the representation of selected literals and state values.
The adversary has a copy of the executable and may use disassemblers,
decompilers, symbolic execution, dynamic instrumentation, and binary
patching. We assume no trusted execution boundary around the ordinary
client process. Code and data available to that process may eventually
be observed by an attacker with equivalent privileges.

Kagura's intended role is to complicate selected analysis tasks while
preserving application behavior. It does not make a client check an
authoritative security decision. It also does not provide a
cryptographic confidentiality guarantee for values that the client
must decrypt. In particular, obscuring a checksum or token-validation
routine does not turn its algorithm into a secure authentication
construction.

\subsection{Claims Examined Here}

We distinguish transformation evidence, behavioral evidence, and resistance
evidence. A changed representation demonstrates a transformation; matching
observations demonstrate agreement on the sampled inputs; resistance requires
an attacker completing a defined task under controlled conditions.
This paper measures code size, static instruction counts for the selected
subset, bounded execution checks, IR repeatability, and startup outcomes.
It does not measure attack completion, analyst time, decompiler quality,
symbolic-execution coverage, or runtime overhead.

Code growth is a deployment cost, not a validated proxy for resistance.
An expression can become longer while remaining easy to simplify.
Successful linking establishes that required symbols resolve, not that an
associated runtime check provides useful protection. A detector signal must
also be distinguished from an attack: legitimate loader behavior can trigger
heuristics. Detector accuracy requires its own deployment study.

\section{System Design}
\label{sec:design}

\subsection{Registration, Policy, and Compiler Integration}

Kagura implements module and function transformations using NPM
interfaces~\cite{llvmNPM}. The shared registry,
\texttt{include/kagura/PassRegistry.def}, defines pass names, enable
flags, constructors, and ordering for the regular automatic pipeline.
It also supplies the named-pass parser and policy-key mappings.
This reduces the risk that a new pass is available through one entry
point but missing from another.

The plugin registers an \texttt{OptimizerLast} extension callback.
For ordinary optimized builds, the selected Kagura passes therefore
run after the standard optimizer pipeline. Module passes precede a
module-to-function adaptor that runs the selected function passes.
JSON policy is loaded before this pipeline is constructed. Infrastructure
passes with special triggers are handled separately.

At the evaluated revision, the registry contains 14 module-pass rows,
16 function-pass rows, and six additional named passes for configuration,
automatic selection, debug information, symbol mapping, vtable protection,
and audit output. Metrics snapshots are wired separately. These counts
include instrumentation and infrastructure, not 36 independent obfuscation
algorithms. The registry also generates numeric tuning mappings and
30 runtime link smoke tests. Shipped profiles are generated from a separate
shared \texttt{Profiles.def}, reducing policy drift between integrations.


Function transformations use shared selection logic for annotations,
protection lists, exclusion lists, and hot-path exclusions. Each pass
also checks its own applicability conditions. Module-wide string
transforms use eligibility checks on globals and their uses. Requesting
a pass does not imply that it transforms every function, or that it
finds any eligible input in a particular translation unit.

The project targets LLVM 17--22, but loadable plugins must match their
host LLVM build; NPM does not provide a cross-version binary ABI.
Plugin loading through \texttt{-fpass-plugin} avoids modifications to
LLVM's source tree. On LLVM 17--21, however, Clang parses plugin-specific
\texttt{-mllvm} options before this plugin is loaded. An explicit
Clang--\texttt{opt}--backend workflow avoids that option-parsing issue.
The project also supplies a statically linked IR tool for environments
where LLVM disables loadable plugins. This paper measures only the
LLVM 21.1.8 native plugin workflow described in Section~\ref{sec:method}.

\subsection{Selected Transformations}

Table~\ref{tab:passes} describes the nine passes examined in this paper.
They are a subset of the toolkit, not a claim about its total pass count.

\begin{table}[t]
\centering\footnotesize
\caption{Selected passes in automatic execution order. A requested pass
may leave the input unchanged when no eligible target exists.}
\label{tab:passes}
\begin{tabularx}{\columnwidth}{@{}llX@{}}
\toprule
Pass & Scope & Transformation \\
\midrule
STR & Module & XOR-encode eligible narrow string globals; emit decryption guards. \\
FLA & Function & Replace eligible control flow with a switch dispatcher. \\
BCF & Function & Introduce opaque predicates and bogus paths. \\
SUB & Function & Substitute integer arithmetic and bitwise expressions. \\
CO & Function & Express eligible integer constants through computed expressions. \\
IBR & Function & Route eligible direct calls to same-module definitions through pointer globals. \\
BBR & Function & Permute eligible non-entry block layout. \\
BBS & Function & Split blocks and introduce connecting branches. \\
MVO & Function & XOR-encode eligible local integer storage on stores and decode on loads. \\
\bottomrule
\end{tabularx}
\end{table}

\paragraph{Control flow.}
FLA assigns distinct nonzero state values to blocks and routes
transitions through a dispatcher. BCF adds predicates whose outcomes
are known by construction. For example, $v\mathbin{|}\neg v$ equals
the all-ones word. Such an identity can be recognized by ordinary
bit-vector reasoning; it is not intrinsically difficult for an SMT
solver. BBR changes block order without changing CFG edges, whereas
BBS introduces new blocks. Backend layout and branch folding can
change how these IR transformations appear in machine code.

\paragraph{Expressions and constants.}
SUB uses equivalent integer expressions, including the following
identities for fixed-width bit vectors with arithmetic modulo $2^w$:
\begin{align}
 a+b &\equiv (a\oplus b)+2(a\mathbin{\&}b),\\
 a-b &\equiv (a\oplus b)-2(\neg a\mathbin{\&}b).
\end{align}
These equations describe values; implementing them in LLVM also
requires attention to poison and overflow flags. CO changes constant
representation. Neither transformation guarantees that a later
optimizer or deobfuscator cannot recover a simpler expression.

\paragraph{Data and calls.}
STR's ordinary path decrypts a string on first guarded use and leaves
the buffer available in plaintext afterward. Its keys and decryption
logic are present in the executable. MVO transforms supported integer
allocations only when their uses can be handled consistently; it
rejects escaping pointers, mismatched access types, and volatile or
atomic accesses. IBR's name is historical: the implementation examined
here rewrites eligible direct \emph{calls}, not arbitrary CFG branches.
External-call resolution is a separate pass outside this subset.

\subsection{Ordering and Semantic Safeguards}

For this subset, the automatic order is
\[
\begin{gathered}
\mathrm{STR}\to\mathrm{FLA}\to\mathrm{BCF}\to\mathrm{SUB}\to\mathrm{CO}\\
\to\mathrm{IBR}\to\mathrm{BBR}\to\mathrm{BBS}\to\mathrm{MVO}.
\end{gathered}
\]
This order follows the registry rather than the order of command-line
flags. It is an implementation choice, not a proof that other orders
are invalid or that this order is optimal. The following interactions
illustrate why composition needs explicit safeguards.

\paragraph{Conditional entry into a flattened function.}
If the original entry has a conditional branch, initializing the
dispatcher with only its first successor changes program behavior.
FLA instead selects the initial state using the original condition:
\begin{lstlisting}[language=C++]
InitVal = B.CreateSelect(CondBr->getCondition(),
    ConstantInt::get(Int32Ty, TrueState),
    ConstantInt::get(Int32Ty, FalseState));
B.CreateStore(InitVal, SwitchVar);
\end{lstlisting}
The implementation declines unsupported entry terminators and
functions with exception-handling constructs. It also checks for
collisions when assigning case values, preventing duplicate switch
labels and accidental use of the reserved default state.

\paragraph{Generated state and memory encoding.}
FLA introduces a dispatch allocation that later passes encounter.
MVO excludes generated storage using metadata and name guards, with
a structural fallback: it detects an allocation whose load feeds a
\texttt{switch}. The fallback does not depend on retaining IR value
names, and conservatively excludes user-written storage with the same
shape as well. This policy keeps the compiler-generated dispatcher
outside MVO's transformation domain. It should not be read as a claim
that dispatch state could never be encoded by a coordinated design.

\paragraph{Eligibility and verification.}
Other guards protect exception-handling blocks, \texttt{musttail}
calls, and unsafe memory uses. These checks reduce the supported input
space to avoid known invalid rewrites. LLVM IR verification detects
structural errors, but cannot establish equivalence of two otherwise
well-formed programs. Execution-based tests remain necessary.

\subsection{Game-Facing Runtime Boundary}

The optional \texttt{Protected<T>} template stores a value using an
XOR-encoded primary field and a related shadow field. Reads decode
both and compare them. The runtime-derived key and both encodings
remain in process memory; the implementation does not rotate the key
on a configurable interval. This mechanism can detect inconsistent
updates to its two copies, but coordinated writes or modification of
the check can bypass it. It is not an authenticated memory primitive.

Unity/IL2CPP-related runtime helpers are separate from the narrow
string pass. The evaluated pipeline does not implement or measure
layout-aware protection of arbitrary \texttt{Il2CppString} objects or
of a complete Unity metadata file. Similarly, the supplied telemetry
callback is a weak no-op until an application supplies an implementation.
The basic-block checksum runtime also ships as an always-passing stub.
Neither is included in the selected-pass experiment. The current profile
experiment below exercises additional runtime dependencies.

The pass named \texttt{tamper} does not provide whole-function machine-code
hash verification. It supplies an IR-opcode hash to
\texttt{kagura\_runtime\_hash\_check}, but the runtime does not compare that
hash with a hash of actual function bytes. It checks selected entry
instruction patterns and suspicious symbol names, while mixing the supplied
hash into a volatile tag. The startup self-check also invokes
platform-specific environment detectors. These are bounded heuristic checks,
not a trusted code-integrity boundary.

Several passes explicitly skip WebAssembly. At the evaluated revision
\texttt{cbab339d5486}, Windows runtime symbols required by anti-debug and
external call indirection were missing. Subsequent work merged as PR~79
provides Windows pass checks and loaded-module symbol resolution. The
measurements here retain the earlier revision and do not evaluate those
Windows additions. Integration availability alone does not establish that
every profile is executable on every listed target.

\section{Selected-Pass Evaluation Method}
\label{sec:method}

\subsection{Workload and Configurations}

The workload is the four \texttt{noinline} C functions in
\texttt{scripts/eval/benchmarks/eval\_functions.c}:
\begin{itemize}
  \item T1, \texttt{eval\_verify\_score}: range, delta, and arithmetic
        consistency checks on a proposed score.
  \item T2, \texttt{eval\_verify\_item}: item and quantity bounds,
        a level restriction, and an FNV-style checksum.
  \item T3, \texttt{eval\_check\_movement}: squared displacement and
        vertical-movement checks.
  \item T4, \texttt{eval\_validate\_token}: a 32-byte accumulator
        computation and comparison with an embedded value.
\end{itemize}
These functions are synthetic verification examples. They are not a
production anti-cheat workload or recommended cryptographic algorithms.

We define \textbf{Plain} as the unmodified baseline, \textbf{Core} as
STR, FLA, BCF, SUB, and CO, and \textbf{Extended} as Core plus IBR,
BBR, BBS, and MVO. These names denote experimental pass lists, not the
project's shipping \texttt{BALANCED} and \texttt{STRONG} JSON presets.
Using the latter would enable different passes and runtime dependencies.

Both transformed configurations use seeds 1, 42, and 12345, BCF
probability 30\%, one BCF iteration, and one SUB iteration. BBS uses a
minimum of three instructions and a maximum of two splits per original
block. We compile Plain once because it has no Kagura random seed.
The six seeded builds describe variability at these chosen settings;
three seeds are not a statistically representative sample of all
possible builds.

\subsection{Build Procedure and Provenance}

Measurements use LLVM 21.1.8 on an arm64 host running macOS 26.7.1.
The compiler target is \texttt{arm64-apple-darwin25.6.0}. The plugin
was built in Release mode from repository revision
\texttt{cbab339d5486}; the artifact records its full revision and
SHA-256 digest. The evaluation files accompany this revision as the
paper artifact, with their own source hashes.

The runner performs four explicit stages:
\begin{enumerate}
  \item Clang emits bitcode with \texttt{-O1} frontend settings and
        \texttt{-Xclang -disable-llvm-passes}.
  \item \texttt{opt} runs \texttt{default<O1>}, loading the plugin for
        Core and Extended. The registered callback appends the selected
        transformations. \texttt{-verify-each} checks IR after passes.
  \item \texttt{llc -O=1 -filetype=obj} generates native code with
        position-independent relocation. This stage avoids running a
        second Clang IR optimization pipeline over the transformed IR.
  \item Clang links that object with the same separately compiled
        \texttt{-O1} self-test harness, without LTO.
\end{enumerate}

This is an explicitly staged experiment, not a claim of byte-for-byte
identity with every Clang driver or application build configuration.
The toolchain versions, commands, source hashes, plugin hash, target,
and per-build results are retained. A repeat transformation in a new
\texttt{opt} process checks exact equality of the emitted IR for each
configuration and seed.

\subsection{Measurements and Acceptance Criteria}

\textbf{Functional checks.} The existing harness has seven assertions:
two score rejection cases, one accepted item and one null item, one
accepted movement and one excessive movement, and one null-token
rejection. Each executable must exit successfully and report all seven
passing. This is a smoke test: it does not exercise the non-null token
loop or establish equivalence for all inputs.

\textbf{Text size.} We read the linked executable's \texttt{\_\_text}
section with \texttt{llvm-objdump}. This includes the unchanged harness
as well as T1--T4. Whole-file byte counts are retained in JSON, but
Mach-O segment layout and metadata make them a less direct measure of
code growth.

\textbf{Instructions.} We count disassembled machine instructions
within each of the four named function symbols, then sum those counts.
We exclude the harness and linker stubs from this metric. Text-size
and instruction-count ratios therefore have different denominators;
they must not be presented as interchangeable measures. Neither count
is a dynamic instruction count or an execution-time measurement.

\section{Selected-Pass Results}
\label{sec:results}

% Generated by run_eval.py; do not edit.
\newcommand{\EvalLLVM}{21.1.8}
\newcommand{\EvalSeeds}{1, 42, 12345}
\begin{table}[t]
\centering\footnotesize
\caption{Native measurements. Ranges span three fixed seeds.
Instructions cover T1--T4; text includes the unchanged harness.}
\label{tab:overhead}
\setlength{\tabcolsep}{3pt}
\begin{tabular}{lrrrr}
\toprule
Config. & Text (B) & Text ratio & Instructions & Instr. ratio \\
\midrule
Plain & 856 & 1.00 & 111 & 1.00 \\
Core & 1,404--1,420 & 1.64--1.66 & 248--252 & 2.23--2.27 \\
Extended & 1,568--1,620 & 1.83--1.89 & 289--302 & 2.60--2.72 \\
\bottomrule
\end{tabular}
\end{table}
\begin{table}[t]
\centering\footnotesize
\caption{Per-function machine instructions: median [min, max]
across seeds. Plain is compiled once.}
\label{tab:perfunc}
\begin{tabular}{lrrr}
\toprule
Function & Plain & Core & Extended \\
\midrule
T1 & 21 & 42 [42, 42] & 46 [46, 46] \\
T2 & 37 & 84 [83, 93] & 98 [97, 105] \\
T3 & 24 & 50 [45, 50] & 54 [49, 54] \\
T4 & 29 & 73 [72, 73] & 97 [96, 100] \\
\bottomrule
\end{tabular}
\end{table}


\subsection{Functional Checks and Repeatability}

Plain and all six transformed executables pass all seven assertions.
Each of the six repeated transformations emits byte-identical IR for
its fixed input and seed. IR verification succeeds throughout the
measured optimizer runs. These results support the limited checks
specified above; they do not establish binary reproducibility across
machines, LLVM releases, input paths, or linkers.

\subsection{Code-Size Costs}

Table~\ref{tab:overhead} reports the measured ranges. Plain has an
856-byte text section. Core occupies 1,404--1,420 bytes
($1.64$--$1.66\times$), and Extended occupies 1,568--1,620 bytes
($1.83$--$1.89\times$). The whole executable remains within
33,800--33,808 bytes for these builds, illustrating why file size alone
would obscure much of this difference.

The four target functions contain 111 instructions in Plain,
248--252 in Core, and 289--302 in Extended. The corresponding ratios
are $2.23$--$2.27\times$ and $2.60$--$2.72\times$.
Table~\ref{tab:perfunc} reports per-function medians and ranges.
The fixed harness contributes to the text denominator but not to these
function counts, explaining why the two aggregate ratios differ.

\subsection{What the Configurations Exercise}

The transformed translation unit contains no eligible string literals
and no eligible calls to other functions defined within that module
after the standard optimizer. The output IR contains no STR decryption
helpers or IBR pointer loads. Thus, enabling STR and IBR in these pass
lists does not provide a measurement of their overhead or effectiveness.
The harness's output string is compiled separately and is not protected.

The emitted IR contains flattening dispatchers, and Extended includes
MVO rewrites. However, aggregate differences between Core and Extended
cannot isolate the marginal cost of BBR, BBS, or MVO. They compose,
and their random-number consumption can affect later choices.
An ablation study with inputs that exercise each pass would be required
to attribute effects to individual transformations.


\section{Current-Profile Deployment Evaluation}
\label{sec:profiles}

\subsection{Settings and Independent Build Protocol}

The previous experiment isolates a selected pass catalogue. We separately
test FAST, BALANCED, and STRONG as shipped, because their policy settings
and runtime dependencies differ from Core and Extended.
FAST enables STR, symbol visibility, and anti-debug instrumentation.
BALANCED enables STR, wide-string protection, BCF, BBR, BBS, dead-code
insertion, global encryption, MVO, symbol visibility, anti-debug, and tamper
instrumentation. STRONG adds AES strings, FLA, SUB, CO, IBR, loop
transformation, and honey values. BCF probability is 30/20/50\% for the
three profiles, with BCF and SUB iteration counts of 1/1/2.
STRONG does not enable every available pass.

This experiment uses the same source revision, LLVM~21.1.8, and Apple M4 Pro
host running macOS~26.7.1, but has an independently recorded build protocol.
Clang emits subject IR with \texttt{-O1 -Xclang -disable-llvm-passes};
\texttt{opt} runs \texttt{default<O1>} with strict JSON policy validation
and \texttt{-verify-each}. Clang then lowers the resulting IR at
\texttt{-O0}, avoiding another ordinary IR optimization pipeline.
Drivers are compiled separately at \texttt{-O1}; protected builds link the
shipped static runtime. We test seeds 7, 42, and 12345. Each transformation
runs twice in separate processes to check IR repeatability.

This differs from the previous experiment's \texttt{llc -O=1} backend
setting. Consequently, the baseline, denominators, and code-size ratios
in this section are independent; differences between the two experiments
cannot be attributed to profile selection alone. The explicit protocol
serves reproducibility, not a claim that these settings are optimal for
production builds.

\subsection{Expanded Differential Inputs}

In addition to the seven existing self-tests, an unobfuscated driver supplies
1,024 deterministic inputs to each of the four subject functions. It covers
accepted and rejected score, item, and movement cases, null pointers,
and non-null token buffers that execute the accumulator loop.
Coordinates and ticks are bounded to avoid signed overflow in the movement
subject. Running the baseline and this driver under LLVM's undefined-behavior
sanitizer reports no violation for the selected sequence.

The baseline accepts 456 score cases, 332 item cases, 480 movement cases,
and zero token cases. Token acceptance is therefore untested. The complete
4,096-return-value output is compared byte for byte with the baseline;
no aggregate digest is used as the comparison oracle. All profiles and seeds
receive the same sequence, so repeated builds are not independent samples
of application behavior. Enabled string and call transformations still lack
eligible targets in this subject module.

\subsection{Compiler Tests and Startup Behavior}

All 84 registered CTest tests pass, including 30 link smoke tests, eight
regression-corpus tests, and one CTest entry containing 48 FileCheck tests.
The 48 nested cases are not additional CTest entries. The existing
differential script reports eight subjects passed with no failures or skips;
its pipeline is a selected subset, not STRONG. The existing fixed-seed
reproducibility script also passes for its four-pass subset.

Nevertheless, all nine shipped profile/seed combinations abort with
\texttt{SIGABRT} before either driver produces its expected output.
A diagnostic executable linked with the same runtime returns one from
\texttt{kagura\_check\_got\_hooks} and zero from the other five probes in
the startup chain. We did not intentionally install hooks or attach a
debugger. The startup chain invokes an aborting response on a positive
signal. This identifies a deployment incompatibility requiring investigation;
it does not measure successful attack detection or a population
false-positive rate.

To separate the environment probes from the remaining transformations,
we repeat all configurations with only
\texttt{passes.anti\_debug=false}. We denote these ablations by a superscript
minus. All nine ablated configurations pass the seven self-tests and match
all 4,096 differential return values. BALANCED$^{-}$ and STRONG$^{-}$
retain tamper instrumentation and the linked runtime. These are successful
ablations, not successful executions of the original shipped profiles.
All 18 protected configurations, including those that abort, produce
byte-identical IR in their paired runs.

% Generated by paper/render_results.py from the selected evaluation JSON.
\begin{table*}[t]
\centering\footnotesize
\caption{Profile observations on arm64 macOS. Ranges span seeds 7, 42, and 12345; plain is built once. A superscript $^{-}$ disables only anti-debug instrumentation. Sizes are bytes; executable sizes use the seven-case self-test driver. Completed runs pass both that driver and the 4,096-value differential comparison.}
\label{tab:current-results}
\begin{tabular}{@{}llrrrrr@{}}
\toprule
Configuration & Execution outcome & Subject text & Text ratio & Executable text & File size & Completed \\
\midrule
Plain & Matches baseline & 880 & $1.00\times$ & 1{,}292 & 33{,}808 & 1/1 \\
FAST & Startup abort & 1{,}076 & $1.22\times$ & 4{,}172 & 36{,}272 & 0/3 \\
BALANCED & Startup abort & 1{,}448--1{,}644 & $1.65\text{--}1.87\times$ & 5{,}776--5{,}972 & 37{,}456--37{,}464 & 0/3 \\
STRONG & Startup abort & 7{,}284--8{,}472 & $8.28\text{--}9.63\times$ & 11{,}612--12{,}800 & 37{,}504--54{,}016 & 0/3 \\
\midrule
FAST$^{-}$ & Matches baseline & 880 & $1.00\times$ & 1{,}292 & 33{,}728 & 3/3 \\
BALANCED$^{-}$ & Matches baseline & 1{,}252--1{,}448 & $1.42\text{--}1.65\times$ & 3{,}768--3{,}964 & 35{,}584--35{,}592 & 3/3 \\
STRONG$^{-}$ & Matches baseline & 7{,}088--8{,}276 & $8.05\text{--}9.40\times$ & 9{,}604--10{,}792 & 35{,}616 & 3/3 \\
\bottomrule
\end{tabular}
\end{table*}


\subsection{Code Growth and Interpretation}

Table~\ref{tab:current-results} distinguishes subject-object text,
linked-executable text, and total file size. The subject object includes
helpers generated within the subject module; the executable also includes
the unchanged self-test driver and linked runtime code. The plain subject
contains 880 text bytes. Successful ablations occupy 880 bytes for FAST$^{-}$,
1,252--1,448 for BALANCED$^{-}$, and 7,088--8,276 for STRONG$^{-}$,
corresponding to $1.00\times$, $1.42$--$1.65\times$, and
$8.05$--$9.40\times$. Ranges cover three fixed seeds and are not confidence
intervals. File sizes additionally reflect layout and alignment.

These observations show that component and link tests can pass while
complete shipped policies fail to execute on a host. Code growth for an
aborting configuration is measurable, but cannot be paired with a claim
of preserved application behavior. The successful ablation provides a
bounded behavioral result for this input set, not evidence of resistance
to an adversary or of acceptable runtime cost.

\section{Limitations and Evaluation Priorities}
\label{sec:limits}

\subsection{Correctness and Workload Coverage}

The seven assertions cover only a small fraction of the workload's
behaviors. In particular, passing a null pointer to T4 returns before
its accumulator loop executes. Section~\ref{sec:profiles} extends this with bounded differential inputs,
including non-null token buffers, but still does not cover accepting token
inputs or arbitrary application behavior. Representative C++ code and
broader input coverage remain necessary.

The original workload also avoids several difficult input shapes.
It cannot establish safety for exception handling, aliasing, concurrent
access, vectors, or arbitrary pointer use. T3 contains signed arithmetic
whose behavior is undefined on some extreme inputs, so an expanded
oracle must either constrain its domain or first define the intended
arithmetic. Obfuscation is not a remedy for undefined source behavior.

\subsection{Security Evidence}

We do not assign analyst-hours or a security ranking to these
configurations. Multiplying assumed per-pass effort factors lacks
empirical calibration and ignores shared simplification steps. The
same recovery technique may remove artifacts introduced by several
passes, and a pass that finds no eligible target adds no protection.
Even a relative ranking requires a defined attack task and validation.

Dynamic tracing and synthesis-based deobfuscation demonstrate why
syntactic complexity alone is insufficient~\cite{yadegari2015generic,
blazytko2017syntia}. XOR keys can be read from the executable or process,
plaintext can be observed during use, and dispatcher state may be
recovered by analysis. None of the reported size ratios quantifies the
effort of these attacks.

A useful resistance experiment would specify a recovery objective,
validate the unobfuscated baseline, use identical harnesses and resource
budgets, and report success, failure, timeout, and infrastructure errors
separately. For symbolic execution it must also define symbolic
arguments, memory initialization, termination conditions, and the
coverage metric. A loader failure or a verdict based only on branch counts cannot establish
resistance. The bundled symbolic-execution script assigns \texttt{PARTIAL}
from a branch threshold, accumulates active states rather than unique
completed paths, and suppresses analysis exceptions. Its verdicts are
therefore excluded from security evidence. No controlled attack comparison
is reported here.

\subsection{Runtime and Platform Costs}

Static instruction counts do not capture branch prediction, cache
behavior, register pressure, startup decryption, or frame-time tails.
We have not measured execution time, energy, compilation time, or
memory consumption. The native macOS experiment also does not establish
overhead on iOS devices, Android devices, x86-64, Windows, or WebAssembly.
Supporting a target in the project is distinct from evaluating each
pass on that target.

The next evaluation should combine representative game workloads with
per-pass ablations and input coverage, then measure runtime costs and
specific deobfuscation tasks. This would allow a deployment decision
based on observed tradeoffs, rather than pass count or code size alone.

\section{Related Work}

Collberg et al. distinguish properties such as potency, resilience,
and execution cost~\cite{collberg1997taxonomy}. That distinction motivates
our separation of transformation structure, code-size overhead, and
unmeasured security outcomes. Control-flow flattening is an established
technique with prior structural analyses~\cite{laszlo2009obfuscating}.

Obfuscator-LLVM demonstrated compiler-integrated protection through
LLVM transformations~\cite{junod2015obfuscator}. O-MVLL is another
open-source LLVM obfuscator, distributed under the Apache 2.0 license,
with Python-driven configuration and mobile toolchain
integration~\cite{omvll}. Tigress provides a broad source-to-source C
transformation framework~\cite{tigress}. Kagura's contribution is an
alternative implementation and integration approach; we do not claim
that open-source plugin deployment is unique to Kagura or that its
security exceeds these tools without a comparative experiment.

Yadegari et al. study automatic deobfuscation of executable
code~\cite{yadegari2015generic}, while Syntia recovers the semantics
of obfuscated code using synthesis~\cite{blazytko2017syntia}.
Banescu et al. evaluate obfuscation against symbolic-execution
attacks~\cite{banescu2016code}. These works motivate evaluations that
measure explicit attacker outcomes rather than infer them from static
code growth. Their findings do not supply calibrated effort multipliers
for Kagura's particular implementation or workload.

\section{Conclusion}

Kagura combines established obfuscation primitives with shared registration,
policy selection, and guards for transformation interactions. The selected-pass
experiment provides a reproducible code-size baseline. The independent
current-profile study finds that passing compiler tests does not ensure
successful deployment: startup checks abort on the evaluation host, while
explicit anti-debug ablations match the bounded differential observations.
These results make configuration, semantic checks, runtime behavior, and
measurement scope visible. Attack resistance, broader platform validation,
and production execution costs remain open evaluation requirements.

\section*{Artifact Availability}

The source repository is \url{https://github.com/qusamulab/kagura}.
Reproduction instructions are in \path{paper/README.md}.
The selected-pass runner writes \path{native_eval.json} and its tables.
The current-profile runner, \path{paper/evaluate.py}, records source and
binary hashes, policies, commands, and individual observations in
\path{paper/results/current_eval.json}. The companion
\path{paper/render_results.py} generates its table. Compiler-test logs
and a runtime diagnostic accompany the artifact. The older exploratory
\path{eval_summary.json} supplies no results in this paper.

\bibliographystyle{IEEEtran}
\newpage
\bibliography{refs}
\end{document}
